% ============================================================
%  LAB REPORT TEMPLATE — Modern Grayscale
%  Compile with: pdflatex lab_report_template.tex
% ============================================================

\documentclass[12pt, a4paper]{report}

% ── Encoding & Language ──────────────────────────────────────
\usepackage[T1]{fontenc}
\usepackage[utf8]{inputenc}
\usepackage[english]{babel}

% ── Modern Typeface (Source Sans / Source Serif) ─────────────
\usepackage{sourcesans}          % sans-serif body (replaces sourcesanspro)
\usepackage{sourceserif}         % serif for display/headings (replaces sourceserifpro)
\renewcommand{\familydefault}{\sfdefault}

% ── Page Geometry ────────────────────────────────────────────
\usepackage[
  top=2.5cm, bottom=2.5cm,
  left=2.8cm, right=2.8cm,
  headheight=14pt
]{geometry}

% ── Color (grayscale palette) ────────────────────────────────
\usepackage{xcolor}
\definecolor{ink}{gray}{0.08}          % near-black text
\definecolor{accent}{gray}{0.22}       % headings / rules
\definecolor{mid}{gray}{0.48}          % captions / labels
\definecolor{light}{gray}{0.88}        % rule fills / table rows
\definecolor{verylightgray}{gray}{0.95}% shaded boxes

% ── Typography & Spacing ─────────────────────────────────────
\usepackage{microtype}
\usepackage{setspace}
\setstretch{1.25}
\usepackage{parskip}
\setlength{\parindent}{0pt}

% ── Headers & Footers ────────────────────────────────────────
\usepackage{fancyhdr}
\pagestyle{fancy}
\fancyhf{}
\renewcommand{\headrulewidth}{0.4pt}
\renewcommand{\headrule}{\hbox to\headwidth{\color{light}\leaders\hrule height \headrulewidth\hfill}}
\fancyhead[L]{\color{mid}\small\textit{Lab \labnum\ — \labtitle}}
\fancyhead[R]{\color{mid}\small\coursecode}
\fancyfoot[C]{\color{mid}\small\thepage}

% ── Section Styling ──────────────────────────────────────────
\usepackage{titlesec}

% Chapter (used as Q-style heading override — see below)
\titleformat{\chapter}[display]
  {\color{accent}\sffamily\Large\bfseries}
  {}
  {0pt}
  {\huge}
\titlespacing*{\chapter}{0pt}{-30pt}{12pt}

% Section → Q-level heading
\titleformat{\section}
  {\color{accent}\sffamily\large\bfseries}
  {}
  {0pt}
  {}
\titlespacing*{\section}{0pt}{18pt}{6pt}

% Subsection
\titleformat{\subsection}
  {\color{accent}\sffamily\normalsize\bfseries}
  {}
  {0pt}
  {}
\titlespacing*{\subsection}{0pt}{12pt}{4pt}

% ── Table of Contents ────────────────────────────────────────
\usepackage{titletoc}
\usepackage{tocloft}
\renewcommand{\cftchapfont}{\sffamily\bfseries\color{accent}}
\renewcommand{\cftsecfont}{\sffamily\color{ink}}
\renewcommand{\cftchapleader}{\color{light}\cftdotfill{\cftdotsep}}
\renewcommand{\cftsecleader}{\color{light}\cftdotfill{\cftdotsep}}
\renewcommand{\cftchappagefont}{\sffamily\bfseries\color{accent}}
\renewcommand{\cftsecpagefont}{\sffamily\color{mid}}
\setlength{\cftbeforechapskip}{6pt}

% ── Figures & Captions ───────────────────────────────────────
\usepackage{graphicx}
\usepackage{float}
\usepackage[
  labelfont={bf, color=accent},
  textfont={color=mid},
  font=small,
  justification=centering,
  skip=6pt
]{caption}

% Figure frame / box
\usepackage{mdframed}
\mdfdefinestyle{figbox}{
  linecolor=ink,
  linewidth=0.4pt,
  innertopmargin=2pt,
  innerbottommargin=2pt,
  innerrightmargin=2pt,
  innerleftmargin=2pt,
  backgroundcolor=white,
  roundcorner=0pt
}

% ── Highlighted Answer Box ───────────────────────────────────
\mdfdefinestyle{answerbox}{
  linecolor=accent,
  linewidth=1.2pt,
  innertopmargin=8pt,
  innerbottommargin=8pt,
  innerrightmargin=10pt,
  innerleftmargin=10pt,
  backgroundcolor=verylightgray,
  roundcorner=4pt
}

% ── Hyperlinks ───────────────────────────────────────────────
\usepackage[
  colorlinks=true,
  linkcolor=accent,
  urlcolor=accent,
  citecolor=mid,
  pdftitle={Lab Report},
  pdfauthor={Team}
]{hyperref}

% ── Misc Utilities ───────────────────────────────────────────
\usepackage{booktabs}
\usepackage{enumitem}
\setlist[itemize]{leftmargin=*, itemsep=2pt, topsep=4pt}
\usepackage{lipsum}     % remove in real report

% ── Custom Commands ──────────────────────────────────────────

% Bold inline answer  →  \ans{value}
\newcommand{\ans}[1]{%
  \textbf{\color{accent}#1}%
}

% Module heading  →  \module{Module Title}
\newcounter{modnum}
\newcommand{\module}[1]{%
  \stepcounter{modnum}%
  \setcounter{qnum}{0}%
  \chapter*{Module \arabic{modnum}: #1}%
  \addcontentsline{toc}{chapter}{Module \arabic{modnum}: #1}%
}

% Q&A question heading  →  \question{Question text}
\newcounter{qnum}
\newcommand{\question}[1]{%
  \stepcounter{qnum}%
  \section*{Q\arabic{modnum}.\arabic{qnum})\quad #1}%
  \addcontentsline{toc}{section}{Q\arabic{modnum}.\arabic{qnum}: #1}%
}

% Screenshot figure  →  \screenshot[width]{file}{Caption text}
\newcommand{\screenshot}[3][\linewidth]{%
  \begin{figure}[H]
    \centering
    \begin{mdframed}[style=figbox]
      \centering
      \includegraphics[width=\dimexpr#1-4pt\relax]{#2}
    \end{mdframed}
    \caption{#3}
  \end{figure}%
}

% Inline answer highlight box  →  \ansbox{Full answer sentence}
\newcommand{\ansbox}[1]{%
    \vspace{8pt}%
    \begin{mdframed}[style=answerbox]
        \textbf{Answer:}\quad #1
    \end{mdframed}%
}

% ── Lab Metadata — EDIT THESE ────────────────────────────────
\newcommand{\coursename}{Digital Forensics}
\newcommand{\coursecode}{CBE322}
\newcommand{\labnum}{5}
\newcommand{\labtitle}{Magnet AXIOM: OS, Encryption and Credentials}
\newcommand{\labdate}{March 22, 2026}

% Team members: \member{Name}{Roll Number}
\newcommand{\teamtable}{%
  \begin{tabular}{ll}
    Kandarp Jindal          & 2023BCY0044 \\
    Akshat Ajit Saraswat      & 2023BCY0048 \\
  \end{tabular}%
}

% Initialize module counter
\setcounter{modnum}{2}

% ============================================================
%  DOCUMENT
% ============================================================
\begin{document}

% ── PAGE 1: TITLE PAGE ───────────────────────────────────────
\begin{titlepage}
  \thispagestyle{empty}
  \pagecolor{white}

  % Top rule
  {\color{accent}\rule{\linewidth}{2.5pt}}
  \vspace{0.5em}

  % Course line
  {\color{mid}\sffamily\large
    \textbf{\coursename} \hfill \textbf{\coursecode}
  }

  \vspace{0.3em}
  {\color{light}\rule{\linewidth}{0.8pt}}

  \vspace{3.5cm}

  % Lab number chip
  \begin{center}
    {\color{accent}\sffamily\fontsize{11}{13}\selectfont
      \textls[120]{LAB \labnum}}\\[6pt]
    {\color{ink}\sffamily\fontsize{26}{32}\selectfont\bfseries \labtitle}
  \end{center}

  \vspace{3cm}

  % Date
  \begin{center}
    {\color{mid}\sffamily\large \labdate}
  \end{center}

  \vspace{2.5cm}

  % Team box
\begin{center}
        {\large \teamtable}
\end{center}

  \vfill

  % Bottom rule
  {\color{light}\rule{\linewidth}{0.8pt}}
  \vspace{0.4em}
  {\color{accent}\rule{\linewidth}{2.5pt}}
\end{titlepage}

% ── PAGE 2: TABLE OF CONTENTS ────────────────────────────────
\newpage
\thispagestyle{empty}

\vspace*{1cm}
{\color{accent}\sffamily\Huge\bfseries Index}\\[4pt]
{\color{accent}\rule{\linewidth}{2pt}}
\vspace{1cm}

{
  \setstretch{1.5}
  \tableofcontents
}
\thispagestyle{empty}

\newpage
\setcounter{page}{1}
\pagestyle{fancy}

% ── PAGE 3+: REPORT CONTENT ────────────────────────────────

\module{OS Part 1}

\question{Locate user accounts and user specific hives.}

Axiom lists User Accounts under Artifacts → Operating System → User Accounts. Here we find users \texttt{ryanJ} and \texttt{RJennings}.
\screenshot{lab5img/01.png}{Axiom Examine v9.11: User Accounts artifact showing two users: ryanJ and RJennings}
For the hives we switch to the File System view in AXIOM and navigate to \texttt{C:\textbackslash Users\textbackslash RJennings} where we find the user-specific registry hives: \texttt{NTUSER.DAT}.
\screenshot{lab5img/02.png}{Axiom Examine v9.11: File System view showing RJennings user folder with NTUSER.DAT hive}
We can see the same for \texttt{ryanJ} under \texttt{C:\textbackslash Users\textbackslash ryanJ}.
\screenshot{lab5img/03.png}{Axiom Examine v9.11: File System view showing ryanJ user folder with NTUSER.DAT hive}

\ansbox{The user accounts are \ans{ryanJ} and \ans{RJennings}.}

\question{Identify the allocated and unallocated area in bytes on the disk.}
In Axiom, we can find partition information under File System → Disk Name → Partitions with details on allocated and unallocated space for each partition. We can see \texttt{Partition 1}, \texttt{Partition 3}, \texttt{Partition 4} and \texttt{Partition 5} with their respective allocated and unallocated sizes in bytes.
\screenshot{lab5img/04.png}{Axiom Examine v9.11: File System view showing Partition 1 details}
\screenshot{lab5img/05.png}{Axiom Examine v9.11: File System view showing Partition 3 details}
\screenshot{lab5img/06.png}{Axiom Examine v9.11: File System view showing Partition 4 details}
\screenshot{lab5img/07.png}{Axiom Examine v9.11: File System view showing Partition 5 details}

\begin{table}[H]
  \centering
  \caption{Partition Unallocated Space Summary}
  \label{tab:partition-unallocated-space}
  \begin{tabular}{lrr}
    \toprule
    \textbf{Partition Name} & \textbf{Unallocated Bytes} & \textbf{Unallocated (GB)} \\
    \midrule
    Partition 1 & 215,408,640 & 0.21 \\
    Partition 3 & 170,017,382,400 & 170.01 \\
    Partition 4 & 491,462,656 & 0.49 \\
    Partition 5 & 552,022,016 & 0.55 \\
    \bottomrule
  \end{tabular}
\end{table}

\ansbox{There is a total of \ans{171.26 GB} of unallocated space. Refer to \autoref{tab:partition-unallocated-space} for specific allocation metrics.}

\question{How can an investigator confirm if the accused has reinstalled the OS trying to destroy potential evidence? How can the investigator get the original OS installation date?}
To determine if the OS has been reinstalled, an investigator should:

\begin{enumerate}
  \item Look for a \texttt{windows.old} folder in the File System view, which indicates a previous OS installation.
  \item Check the installation date in System Information and compare it with other file creation dates to identify timeline inconsistencies suggesting a recent OS reinstallation.
\end{enumerate}

The \textbf{OS installation date} can be obtained directly from Artifacts → Operating System → System Information, which displays the exact timestamp of when Windows was installed on the system.
\screenshot{lab5img/08.png}{Axiom Examine v9.11: System Information artifact showing OS installation date as 12-01-2022 13:54:15}

In our case no \texttt{windows.old} folder is present.

\ansbox{Signs of OS reinstallation include: (1) presence of \ans{windows.old} folder and, (2) recent OS installation date compared to user file dates. The OS installation date is \ans{12-01-2022 13:54:15}.}

\question{What is the timezone information of the imaged system?}

Timezone information can be found in Artifacts → Operating System → Timezone Information.
\screenshot{lab5img/09.png}{Axiom Examine v9.11: System Information artifact showing Time Zone as UTC-05:00 Indiana (East)}

\ansbox{The timezone of the imaged system is \ans{UTC-05:00 Indiana (East)}.}

\module{Encryption and Credentials}

\question{Identify the port used for exfiltration}

Apply the display filter \texttt{ip.src == 191.19.11.11} to isolate packets originating from the threat actor.
Inspect the \textbf{Transport Layer} details in the packet dissection panel.

\screenshot{example-image-b}{Wireshark — Filtered view showing repeated SYN packets on port 4444}

The repeated SYN packets confirm a persistent connection attempt on an uncommon port.

\ansbox{The exfiltration port is \ans{4444} (non-standard, typically associated with reverse shells).}

\question{Determine the protocol used in the attack}

Follow the TCP stream (\textbf{right-click → Follow → TCP Stream}) on one of the suspicious packets.
The stream data reveals cleartext commands consistent with a Netcat-style reverse shell session.

\screenshot{example-image-c}{Wireshark — TCP Stream reconstruction showing shell commands in plaintext}

The lack of encryption and the command-response pattern confirm the protocol in use.

\ansbox{The attack leveraged a raw \ans{TCP reverse shell} (Netcat) over port 4444.}

% ──────────────────────────────────────────────────────────────
%  ADD MORE MODULES / QUESTIONS BELOW
%  New module  →  \module{Title}        (resets question counter)
%  New question →  \question{Title}
% ──────────────────────────────────────────────────────────────

\end{document}
